% Diapositiva 14

% Luis
%============================================
% 12. Explicar ese valor mínimo y máximo de la tabla 
%=============================================

\begin{frame}{Resultados: Distancias Calculadas}

    \begin{columns}[c]

        %========================================
        % COLUMNA IZQUIERDA: Tabla de Resultados
        %========================================
        \column{0.58\textwidth}
        \begin{block}{Distancia a cada Sensor ($R_i$ [km])}
            \centering
            \scriptsize 
            \begin{tabular}{@{}lcccc@{}}
                \toprule
            \textbf{Sensor} & \textbf{$x_i$ [km]} & \textbf{$y_i$ [km]} & \textbf{$z_i$ [km]} & \textbf{$R_i$ [km]} \\ \midrule
                S1  & -4.0 & -3.0 &  0.2 & 5.7940 \\
                S2  & -2.5 &  2.8 & -0.1 & 5.2583 \\
                S3  &  0.0 &  4.0 &  0.1 & 5.0912 \\
                S4  &  3.5 &  3.0 & -0.2 & 4.5321 \\
                S5  &  4.5 & -1.5 &  0.0 & 3.5482 \\
                S6  &  2.0 & -4.0 &  0.3 & 3.5833 \\
                S7  & -1.0 & -4.5 & -0.1 & 4.4193 \\
                S8  & -4.5 &  1.0 &  0.2 & 6.1172 \\
                S9  &  0.5 &  0.5 &  0.0 & 1.8412 \\
                S10 &  2.8 &  0.8 & -0.3 & 2.4000 \\
                S11 & -1.8 & -1.2 &  0.1 & 3.2558 \\
                S12 &  1.0 &  3.2 &  0.2 & 4.2107 \\ \bottomrule
            \end{tabular}
        \end{block}

        %========================================
        % COLUMNA DERECHA: Análisis Crítico
        %========================================
        \column{0.38\textwidth}

        \begin{block}{Análisis de la Red}
            \scriptsize
            Resultados tras aplicar la métrica euclidiana desde la fuente real a cada nodo.
        \end{block}

        \vspace{0.4cm}

        \begin{block}{Valores Críticos}
            \scriptsize
            \begin{itemize}
                \item \textbf{Mínimo:} S9 ($R_9 \approx 1.84$ km) \\
                \vspace{0.1cm}
                $\rightarrow$ \textit{Mayor amplitud esperada.}
                
                \vspace{0.4cm}
                
                \item \textbf{Máximo:} S8 ($R_8 \approx 6.12$ km) \\
                \vspace{0.1cm}
                $\rightarrow$ \textit{Mayor atenuación (pérdida).}
            \end{itemize}
        \end{block}

    \end{columns}

\end{frame}